\documentclass[11pt]{article}

\usepackage[utf8]{inputenc}
\usepackage[T1]{fontenc}
\usepackage{lmodern}
\usepackage{geometry}
\usepackage{amsmath,amssymb,amsfonts,amsthm,mathtools}
\usepackage{bm}
\usepackage{enumitem}
\usepackage{microtype}
\usepackage{hyperref}
\usepackage{titlesec}
\usepackage{booktabs}
\usepackage{array}
\usepackage{setspace}
\usepackage{mathrsfs}
\usepackage{physics}

\geometry{margin=1in}
\hypersetup{colorlinks=true, linkcolor=black, urlcolor=blue, citecolor=black}
\setlist[enumerate]{topsep=0.4em,itemsep=0.35em}
\setlist[itemize]{topsep=0.3em,itemsep=0.25em}
\titleformat{\section}{\large\bfseries}{\thesection.}{0.5em}{}
\titleformat{\subsection}{\normalsize\bfseries}{\thesubsection.}{0.5em}{}
\setstretch{1.06}

\newcommand{\Aether}{\AE{}ther}
\newcommand{\Aflow}{\AE{}ther-flow}
\newcommand{\TheoryName}{\Aflow{} Interpretation of Relativity}
\newcommand{\TheoryProgram}{\Aether{} / \Aflow{} framework}
\newcommand{\RespShapeReadout}{\widetilde q_{\mathrm{br}}}
\newcommand{\RespShapeCov}{\widetilde\kappa_{\mathrm{br}}}
\newcommand{\RespShapeRelax}{\widetilde\rho_{\mathrm{br}}}
\newcommand{\RespShapeWell}{\widetilde\lambda_{\mathrm{br}}}
\newcommand{\RespShapeEq}{\widetilde U_{\mathrm{br}}^{\ast}}
\newcommand{\RespShapeIso}{\widetilde\eta_{\mathrm{br}}}
\newcommand{\RespRawReadout}{\widehat q_{\mathrm{br}}}
\newcommand{\RespRawRef}{\widehat U_{\mathrm{br}}^{\mathrm{ref}}}
\newcommand{\RespRawCov}{\widehat\kappa_{\mathrm{br}}}
\newcommand{\RespRawRelax}{\widehat\rho_{\mathrm{br}}}
\newcommand{\RespRawWell}{\widehat\lambda_{\mathrm{br}}}
\newcommand{\RespRawEq}{\widehat U_{\mathrm{br}}^{\ast}}
\newcommand{\RespVacPhase}{\phi_{\mathrm{br}}}
\newcommand{\CurvProfileCan}{F_{\mathrm{br}}^{\mathrm{can}}}
\newcommand{\DownPkg}{\mathscr P_{\mathrm{down}}}
\newcommand{\ShapeTuple}{\mathfrak S_{\mathrm{br}}}
\newcommand{\ScaleMod}{s_{\mathrm{br}}}
\newcommand{\PrimClass}{\mathfrak C_{\mathrm{br}}}
\newcommand{\PrimRead}{r_{\mathrm{br}}}
\newcommand{\PrimCovSlot}{a_{\mathrm{br}}^{\varphi}}
\newcommand{\PrimRelaxSlot}{a_{\mathrm{br}}^{V}}
\newcommand{\PrimKinetic}{\bm a_{\mathrm{br}}}
\newcommand{\PrimWellSlot}{b_{\mathrm{br}}^{\lambda}}
\newcommand{\PrimEqSlot}{b_{\mathrm{br}}^{U}}
\newcommand{\PrimAnchor}{\bm b_{\mathrm{br}}}
\newcommand{\PrimMap}{\Pi_{\mathrm{shape}}}
\newcommand{\RawMap}{\Pi_{\mathrm{raw}}}
\newcommand{\LiftSector}{\mathfrak D_{\mathrm{br}}^{\mathrm{amp}}}
\newcommand{\DeepReadAmp}{\xi_{\mathrm{br}}^{q}}
\newcommand{\DeepCovAmp}{\xi_{\mathrm{br}}^{\varphi}}
\newcommand{\DeepRelaxAmp}{\xi_{\mathrm{br}}^{V}}
\newcommand{\DeepKinAmp}{\bm\xi_{\mathrm{br}}}
\newcommand{\DeepWellAmp}{\zeta_{\mathrm{br}}^{\lambda}}
\newcommand{\DeepEqAmp}{\zeta_{\mathrm{br}}^{U}}
\newcommand{\DeepAnchorAmp}{\bm\zeta_{\mathrm{br}}}
\newcommand{\DeepScaleAmp}{\varsigma_{\mathrm{br}}}
\newcommand{\AmpMap}{\Pi_{\mathrm{amp}}}
\newcommand{\SignQuot}{\sim_{\pm}}
\newcommand{\Rstar}{\mathbb R^{\times}}
\newcommand{\AbsAmpSector}{\mathfrak M_{\mathrm{br}}}
\newcommand{\PolaritySector}{\mathfrak P_{\mathrm{br}}}
\newcommand{\OriginSector}{\mathfrak D_{\mathrm{br}}^{\mathrm{ori}}}
\newcommand{\AbsReadAmp}{\chi_{\mathrm{br}}^{q}}
\newcommand{\AbsCovAmp}{\chi_{\mathrm{br}}^{\varphi}}
\newcommand{\AbsRelaxAmp}{\chi_{\mathrm{br}}^{V}}
\newcommand{\AbsKinAmp}{\bm\chi_{\mathrm{br}}}
\newcommand{\AbsWellAmp}{\upsilon_{\mathrm{br}}^{\lambda}}
\newcommand{\AbsEqAmp}{\upsilon_{\mathrm{br}}^{U}}
\newcommand{\AbsAnchorAmp}{\bm\upsilon_{\mathrm{br}}}
\newcommand{\AbsScaleAmp}{\omega_{\mathrm{br}}}
\newcommand{\PolRead}{\epsilon_{\mathrm{br}}^{q}}
\newcommand{\PolCov}{\epsilon_{\mathrm{br}}^{\varphi}}
\newcommand{\PolRelax}{\epsilon_{\mathrm{br}}^{V}}
\newcommand{\PolWell}{\delta_{\mathrm{br}}^{\lambda}}
\newcommand{\PolEq}{\delta_{\mathrm{br}}^{U}}
\newcommand{\PolScale}{\sigma_{\mathrm{br}}}
\newcommand{\PolTuple}{\bm\epsilon_{\mathrm{br}}}
\newcommand{\OriginMap}{\Pi_{\mathrm{ori}}}
\newcommand{\AbsMap}{\Pi_{\mathrm{abs}}}
\newcommand{\DeepMap}{\Pi_{\mathrm{deep}}}
\newcommand{\PolarityQuot}{\sim_{\mathrm{pol}}}

\newtheorem{definition}{Definition}
\newtheorem{proposition}{Proposition}
\newtheorem{remark}{Remark}
\newtheorem{corollary}{Corollary}
\newtheorem{theorem}{Theorem}

\title{\textbf{The \TheoryName{}}\\[0.4em]
\large Nonlocal Bridge Self-Response Off-Shell Profile-Selection Bridge-Order Orbit-Shape/Modulus Primitive-Class Amplitude-Sector Origin Note}
\author{}
\date{}

\begin{document}

\maketitle

\begin{abstract}
The benchmark exact-theory statement of \TheoryName{} remains fixed and is not reopened here. The overview-first exact-closure package is still the public front door. The present note acts only on the optional deeper derivational bridge, and it begins from the now-frozen primitive-class setup note together with the primitive-class amplitude-lift origin note immediately below it.

Those two notes already fixed the live architecture sharply. The setup note isolated the readout block, the ordered kinetic pair, the radial anchor pair, and the positive orbit modulus as the unique minimal interface below the orbit-shape/modulus burden. The amplitude-lift note then chose the first concrete deeper realization through that interface: the signed amplitude sector
\[
\LiftSector
=
\Rstar\times(\Rstar)^2\times(\Rstar)^2\times\Rstar
\]
with coordinates
\[
(\DeepReadAmp,\DeepKinAmp,\DeepAnchorAmp,\DeepScaleAmp)
=
(\DeepReadAmp,(\DeepCovAmp,\DeepRelaxAmp),(\DeepWellAmp,\DeepEqAmp),\DeepScaleAmp),
\]
whose componentwise square reproduces the primitive class and hence the same raw-orbit lift and the same downstream package. One narrower gap remains. Those six signed amplitudes are still the present deepest recorded layer rather than themselves being generated from still deeper substrate data.

This note supplies the smallest honest continuation below that layer. The new deeper sector is the oriented absolute-amplitude sector
\[
\OriginSector
=
\AbsAmpSector\times\PolaritySector
=
\Bigl(\mathbb R_{>0}\times(\mathbb R_{>0})^2\times(\mathbb R_{>0})^2\times\mathbb R_{>0}\Bigr)
\times
\Bigl(\{\pm1\}\times\{\pm1\}^2\times\{\pm1\}^2\times\{\pm1\}\Bigr),
\]
with coordinates
\[
(\AbsReadAmp,\AbsKinAmp,\AbsAnchorAmp,\AbsScaleAmp;\PolRead,\PolCov,\PolRelax,\PolWell,\PolEq,\PolScale).
\]
The origin map is the componentwise oriented product
\[
\OriginMap:
(\AbsReadAmp,\AbsKinAmp,\AbsAnchorAmp,\AbsScaleAmp;\PolRead,\PolCov,\PolRelax,\PolWell,\PolEq,\PolScale)
\longmapsto
(\PolRead\AbsReadAmp,\PolCov\AbsCovAmp,\PolRelax\AbsRelaxAmp,\PolWell\AbsWellAmp,\PolEq\AbsEqAmp,\PolScale\AbsScaleAmp).
\]
It is a bijection onto the signed amplitude sector, so each signed amplitude factors uniquely into a positive absolute amplitude and a discrete polarity label.

Using the frozen amplitude lift exactly as already fixed, the note then proves four facts. First, the primitive class, the readout block, the ordered kinetic pair, the radial anchor pair, and the positive orbit modulus depend only on the positive absolute-amplitude block. Second, the same raw-orbit lift is reproduced, now written as squares of composite absolute amplitudes. Third, the same downstream density/current block, the same primitive bridge-order block, the same phase-neutral minimizing branch, the same canonical profile, and the same dressed bridge map are recovered unchanged. Fourth, the only quotient structure present is still the finite polarity-forgetting quotient corresponding exactly to the old sign quotient of square-root representatives. It does not identify distinct positive moduli, and the new sector introduces no vacuum-angle locking or angle quotient. Therefore the current verdicts remain unchanged: the positive scale orbit is still a canonical-normalization orbit rather than a proved redundancy, and the global \(O(2)\) vacuum angle remains residual.

The scope is deliberately narrow. This note does not claim a unique final substrate microphysics. Its narrower positive result is that the signed amplitude layer is no longer primitive either: it is the image of an explicit deeper oriented absolute-amplitude sector through a frozen map that preserves the already-positive downstream package exactly.
\end{abstract}

\section{Introduction}

The active repository no longer needs another bridge note in order to justify its public theory claim. The exact-closure sequence already presents \TheoryName{} as the disciplined exact relativistic theory statement built on the \Aether{} / \Aflow{} ontology, and that benchmark package remains the front-door reading order. The present note therefore does not strengthen the flagship theory claim. It acts only on the optional deeper derivational continuation beneath that fixed benchmark.

On that optional continuation line the latest burden is now sharply localized. The orbit-shape/modulus-origin note fixed the deeper architectural task below the raw positive vacuum orbit. The primitive-class setup note then froze the unique minimal interface into that task. The primitive-class amplitude-lift note next realized that interface by choosing the signed amplitude sector \((\DeepReadAmp,\DeepCovAmp,\DeepRelaxAmp,\DeepWellAmp,\DeepEqAmp,\DeepScaleAmp)\) as the first direct lower layer and by proving that the only quotient at that level is the finite sign quotient of square-root representatives. What remains open is narrower again: whether those signed amplitudes themselves can be generated from still deeper substrate structure without disturbing the now-frozen downstream package.

That is the exact burden of the present manuscript. It does not reopen another same-layer orbit/modulus discussion, it does not alter the already-positive raw-orbit lift, and it does not reopen the current scale-orbit or residual-angle verdicts unless the new deeper sector actually forces that revision. Instead it chooses the smallest still-deeper sector compatible with the amplitude-lift result, writes the explicit map from that sector into the six signed amplitudes, and then checks that the already-frozen maps into the primitive class, the orbit-shape data, the raw orbit, and the downstream package remain unchanged.

\paragraph{Framework and claim boundary.}
\TheoryProgram{} denotes the broader \Aether{} / \Aflow{} framework built on the claims that the \Aether{} is the underlying four-dimensional substrate of reality, the \Aflow{} is its intrinsic ordered motion, observed three-dimensional space is the local experiential slice of that deeper substrate, S-time is the experienced order of change arising from matter, light, and the \Aflow{}, and observed expansion is the three-dimensional appearance of deeper four-dimensional ordered motion. Within that framework, \TheoryName{} denotes the exact-closure theory in which the effective gravitational dynamics are adopted to be exactly Einsteinian with universal matter coupling. In the active sequence, \emph{adoption} means use of that established relativistic dynamics without claiming substrate derivation, while \emph{derivation} is reserved for a first-principles recovery from explicit substrate variables.

This note stays strictly inside that boundary. Its positive claim is not that the final unique substrate completion has already been found. Its positive claim is narrower: the six signed amplitudes feeding the frozen primitive class are now themselves realized as oriented readouts of a still deeper sector, and the already-positive raw-orbit and downstream package remain exactly the same.

\section{Frozen Amplitude-Lift Layer}

\begin{definition}[Frozen amplitude sector and primitive-class lift]
\label{def:frozenamp}
The signed amplitude sector fixed by the amplitude-lift note is
\begin{equation}
\LiftSector
:=
\Rstar\times(\Rstar)^2\times(\Rstar)^2\times\Rstar,
\label{eq:liftsector}
\end{equation}
with coordinates
\begin{equation}
(\DeepReadAmp,\DeepKinAmp,\DeepAnchorAmp,\DeepScaleAmp)
=
(\DeepReadAmp,(\DeepCovAmp,\DeepRelaxAmp),(\DeepWellAmp,\DeepEqAmp),\DeepScaleAmp).
\label{eq:liftcoords}
\end{equation}
Its frozen map into the primitive class is
\begin{equation}
\AmpMap:\LiftSector\longrightarrow\PrimClass
\label{eq:ampmap}
\end{equation}
given by
\begin{equation}
\AmpMap(\DeepReadAmp,\DeepKinAmp,\DeepAnchorAmp,\DeepScaleAmp)
=
\bigl(
(\DeepReadAmp)^2,
((\DeepCovAmp)^2,(\DeepRelaxAmp)^2),
((\DeepWellAmp)^2,(\DeepEqAmp)^2),
(\DeepScaleAmp)^2
\bigr).
\label{eq:ampmapexplicit}
\end{equation}
\end{definition}

\begin{definition}[Frozen setup map and raw lift]
\label{def:frozenraw}
The primitive-class setup map remains
\begin{equation}
\PrimMap(\PrimRead,\PrimKinetic,\PrimAnchor,\ScaleMod)
=
(\RespShapeReadout,\RespShapeCov,\RespShapeRelax,\RespShapeWell,\RespShapeEq,\ScaleMod)
\label{eq:primmap}
\end{equation}
with
\begin{equation}
\RespShapeReadout=\PrimRead,
\qquad
\RespShapeCov=\PrimCovSlot,
\qquad
\RespShapeRelax=\PrimRelaxSlot,
\label{eq:shapeA}
\end{equation}
\begin{equation}
\RespShapeWell=\PrimWellSlot,
\qquad
\RespShapeEq=\PrimEqSlot,
\label{eq:shapeB}
\end{equation}
and the frozen raw-orbit lift remains
\begin{equation}
\RespRawReadout=\ScaleMod\,\RespShapeReadout,
\qquad
\RespRawRef=\ScaleMod,
\qquad
\RespRawCov=\ScaleMod^{2}\RespShapeCov,
\label{eq:rawliftA}
\end{equation}
\begin{equation}
\RespRawRelax=\ScaleMod^{2}\RespShapeRelax,
\qquad
\RespRawWell=\RespShapeWell,
\qquad
\RespRawEq=\ScaleMod\,\RespShapeEq.
\label{eq:rawliftB}
\end{equation}
\end{definition}

\begin{definition}[Frozen downstream package]
\label{def:downstream}
At the present derivational stage, the downstream package \(\DownPkg\) means the same already-recorded density/current block, the same primitive bridge-order block, the same phase-neutral minimizing branch, the same canonical profile
\begin{equation}
\CurvProfileCan(x)=1+\RespShapeEq\,x^2
=
1+\frac{\RespRawEq}{\RespRawRef}x^2,
\label{eq:profile}
\end{equation}
and the same dressed bridge map, all understood to depend on deeper data only through the induced pair \((\ShapeTuple,\ScaleMod)\), the raw lift \eqref{eq:rawliftA}--\eqref{eq:rawliftB}, and the same minimizing branch.
\end{definition}

\begin{remark}
Definitions \ref{def:frozenamp}--\ref{def:downstream} are frozen input in this note. The present manuscript does not alter the primitive-class interface or the amplitude-lift formulas. It acts only below them.
\end{remark}

\section{The Deeper Oriented Absolute-Amplitude Sector}

\begin{definition}[Absolute-amplitude block and polarity block]
\label{def:originsector}
Define the positive absolute-amplitude block
\begin{equation}
\AbsAmpSector
:=
\mathbb R_{>0}\times(\mathbb R_{>0})^2\times(\mathbb R_{>0})^2\times\mathbb R_{>0},
\label{eq:abssector}
\end{equation}
with coordinates
\begin{equation}
(\AbsReadAmp,\AbsKinAmp,\AbsAnchorAmp,\AbsScaleAmp)
=
(\AbsReadAmp,(\AbsCovAmp,\AbsRelaxAmp),(\AbsWellAmp,\AbsEqAmp),\AbsScaleAmp),
\label{eq:abscoords}
\end{equation}
and the discrete polarity block
\begin{equation}
\PolaritySector
:=
\{\pm1\}\times\{\pm1\}^2\times\{\pm1\}^2\times\{\pm1\},
\label{eq:polsector}
\end{equation}
with coordinates
\begin{equation}
(\PolRead,\PolCov,\PolRelax,\PolWell,\PolEq,\PolScale).
\label{eq:polcoords}
\end{equation}
The deeper amplitude-sector origin chosen in this note is their product
\begin{equation}
\OriginSector:=\AbsAmpSector\times\PolaritySector.
\label{eq:originsector}
\end{equation}
\end{definition}

\begin{definition}[Origin map into the signed amplitudes]
\label{def:originmap}
Define the origin map
\begin{equation}
\OriginMap:\OriginSector\longrightarrow\LiftSector
\label{eq:originmap}
\end{equation}
by the componentwise oriented product
\begin{align}
\OriginMap(\AbsReadAmp,\AbsKinAmp,\AbsAnchorAmp,\AbsScaleAmp;\PolRead,\PolCov,\PolRelax,\PolWell,\PolEq,\PolScale)
=\;&
(\PolRead\AbsReadAmp,\PolCov\AbsCovAmp,\PolRelax\AbsRelaxAmp,
\nonumber\\
&\PolWell\AbsWellAmp,\PolEq\AbsEqAmp,\PolScale\AbsScaleAmp).
\label{eq:originmapexplicit}
\end{align}
\end{definition}

\begin{proposition}[Unique factorization of the signed amplitudes]
\label{prop:bijection}
The origin map \(\OriginMap\) is a bijection. Equivalently, every signed amplitude tuple in \(\LiftSector\) factors uniquely into a positive absolute-amplitude block and a discrete polarity block.
\end{proposition}

\begin{proof}
Given any point of \(\OriginSector\), every component of \eqref{eq:originmapexplicit} is nonzero because the absolute-amplitude coordinates are positive and the polarity labels are \(\pm1\). Hence \(\OriginMap\) lands in \(\LiftSector\).

Conversely, given any
\[
(\DeepReadAmp,\DeepCovAmp,\DeepRelaxAmp,\DeepWellAmp,\DeepEqAmp,\DeepScaleAmp)\in\LiftSector,
\]
define
\[
\AbsReadAmp=\abs{\DeepReadAmp},
\quad
\AbsCovAmp=\abs{\DeepCovAmp},
\quad
\AbsRelaxAmp=\abs{\DeepRelaxAmp},
\quad
\AbsWellAmp=\abs{\DeepWellAmp},
\quad
\AbsEqAmp=\abs{\DeepEqAmp},
\quad
\AbsScaleAmp=\abs{\DeepScaleAmp},
\]
and
\[
\PolRead=\frac{\DeepReadAmp}{\abs{\DeepReadAmp}},
\quad
\PolCov=\frac{\DeepCovAmp}{\abs{\DeepCovAmp}},
\quad
\PolRelax=\frac{\DeepRelaxAmp}{\abs{\DeepRelaxAmp}},
\quad
\PolWell=\frac{\DeepWellAmp}{\abs{\DeepWellAmp}},
\quad
\PolEq=\frac{\DeepEqAmp}{\abs{\DeepEqAmp}},
\quad
\PolScale=\frac{\DeepScaleAmp}{\abs{\DeepScaleAmp}}.
\]
Then \(\OriginMap\) returns the original signed amplitude tuple. This proves surjectivity. The same formulas also give uniqueness, because the absolute value and the sign of each nonzero real number are unique.
\end{proof}

\begin{corollary}[The amplitude layer is no longer primitive]
\label{cor:notprimitive}
At the level of the present note, the signed amplitude layer \(\LiftSector\) is not primitive. It is the image of the oriented absolute-amplitude sector \(\OriginSector\) under the bijection \(\OriginMap\).
\end{corollary}

\begin{proof}
This is exactly Proposition \ref{prop:bijection}.
\end{proof}

\begin{remark}
Corollary \ref{cor:notprimitive} is the narrow positive content of the present continuation. The note does not claim a unique microscopic completion beyond \(\OriginSector\). It claims only that the signed amplitude variables themselves are no longer terminal input.
\end{remark}

\section{Map Into the Primitive Class and Orbit-Shape Data}

\begin{definition}[Absolute-amplitude map into the primitive class]
\label{def:absmap}
Define the positive map
\begin{equation}
\AbsMap:\OriginSector\longrightarrow\PrimClass
\label{eq:absmap}
\end{equation}
by
\begin{equation}
\AbsMap:=\AmpMap\circ\OriginMap.
\label{eq:absmapdef}
\end{equation}
Explicitly,
\begin{align}
\AbsMap
=\;&
\bigl(
(\AbsReadAmp)^2,
((\AbsCovAmp)^2,(\AbsRelaxAmp)^2),
((\AbsWellAmp)^2,(\AbsEqAmp)^2),
(\AbsScaleAmp)^2
\bigr).
\label{eq:absexplicit}
\end{align}
\end{definition}

\begin{proposition}[The primitive class depends only on the positive absolute-amplitude block]
\label{prop:absonly}
The map \(\AbsMap\) is independent of the polarity coordinates
\((\PolRead,\PolCov,\PolRelax,\PolWell,\PolEq,\PolScale)\). In particular,
\begin{equation}
\PrimRead=(\AbsReadAmp)^2,
\qquad
\PrimCovSlot=(\AbsCovAmp)^2,
\qquad
\PrimRelaxSlot=(\AbsRelaxAmp)^2,
\label{eq:primA}
\end{equation}
\begin{equation}
\PrimWellSlot=(\AbsWellAmp)^2,
\qquad
\PrimEqSlot=(\AbsEqAmp)^2,
\qquad
\ScaleMod=(\AbsScaleAmp)^2.
\label{eq:primB}
\end{equation}
\end{proposition}

\begin{proof}
Substitute \eqref{eq:originmapexplicit} into the componentwise squares \eqref{eq:ampmapexplicit}. Every polarity coordinate appears only squared, hence drops out. The resulting primitive-class formulas are exactly \eqref{eq:primA}--\eqref{eq:primB}.
\end{proof}

\begin{definition}[Composite deeper map into the orbit-shape/modulus data]
\label{def:deepmap}
Define
\begin{equation}
\DeepMap:=\PrimMap\circ\AbsMap:\OriginSector\longrightarrow(\ShapeTuple,\ScaleMod).
\label{eq:deepmap}
\end{equation}
Explicitly,
\begin{equation}
\RespShapeReadout=(\AbsReadAmp)^2,
\qquad
\RespShapeCov=(\AbsCovAmp)^2,
\qquad
\RespShapeRelax=(\AbsRelaxAmp)^2,
\label{eq:deepshapeA}
\end{equation}
\begin{equation}
\RespShapeWell=(\AbsWellAmp)^2,
\qquad
\RespShapeEq=(\AbsEqAmp)^2,
\qquad
\ScaleMod=(\AbsScaleAmp)^2.
\label{eq:deepshapeB}
\end{equation}
\end{definition}

\begin{corollary}[Map into the readout block, ordered kinetic pair, radial anchor pair, and positive orbit modulus]
\label{cor:blockmap}
Under the deeper map \(\DeepMap\), the frozen interface blocks are recovered as
\begin{enumerate}
    \item the readout block
    \[
    \RespShapeReadout=(\AbsReadAmp)^2;
    \]
    \item the ordered kinetic pair
    \[
    (\RespShapeCov,\RespShapeRelax)=((\AbsCovAmp)^2,(\AbsRelaxAmp)^2);
    \]
    \item the radial anchor pair
    \[
    (\RespShapeWell,\RespShapeEq)=((\AbsWellAmp)^2,(\AbsEqAmp)^2);
    \]
    \item the positive orbit modulus
    \[
    \ScaleMod=(\AbsScaleAmp)^2.
    \]
\end{enumerate}
Thus the new deeper sector factors through the already-frozen interface exactly as required.
\end{corollary}

\begin{proof}
This is just the explicit content of Definition \ref{def:deepmap}.
\end{proof}

\begin{remark}
Corollary \ref{cor:blockmap} shows that the present note is not another same-layer interface-fixing manuscript. The interface was already frozen. What changes here is only the deeper origin of the six signed amplitudes feeding that interface.
\end{remark}

\section{Same Raw-Orbit Lift and Same Downstream Package}

\begin{theorem}[The raw-orbit lift is unchanged]
\label{thm:rawlift}
Composing \(\DeepMap\) with the frozen raw lift \eqref{eq:rawliftA}--\eqref{eq:rawliftB} yields exactly the same raw orbit as in the amplitude-lift note. In the present variables,
\begin{equation}
\RespRawReadout=(\AbsScaleAmp\,\AbsReadAmp)^2,
\qquad
\RespRawRef=(\AbsScaleAmp)^2,
\qquad
\RespRawCov=(\AbsScaleAmp\,\AbsCovAmp)^2,
\label{eq:rawabsA}
\end{equation}
\begin{equation}
\RespRawRelax=(\AbsScaleAmp\,\AbsRelaxAmp)^2,
\qquad
\RespRawWell=(\AbsWellAmp)^2,
\qquad
\RespRawEq=(\AbsScaleAmp\,\AbsEqAmp)^2.
\label{eq:rawabsB}
\end{equation}
\end{theorem}

\begin{proof}
Insert \eqref{eq:deepshapeA}--\eqref{eq:deepshapeB} into the frozen raw lift \eqref{eq:rawliftA}--\eqref{eq:rawliftB}. This gives
\[
\RespRawReadout=(\AbsScaleAmp)^2(\AbsReadAmp)^2=(\AbsScaleAmp\,\AbsReadAmp)^2,
\qquad
\RespRawRef=(\AbsScaleAmp)^2,
\]
\[
\RespRawCov=(\AbsScaleAmp)^2(\AbsCovAmp)^2=(\AbsScaleAmp\,\AbsCovAmp)^2,
\qquad
\RespRawRelax=(\AbsScaleAmp)^2(\AbsRelaxAmp)^2=(\AbsScaleAmp\,\AbsRelaxAmp)^2,
\]
\[
\RespRawWell=(\AbsWellAmp)^2,
\qquad
\RespRawEq=(\AbsScaleAmp)^2(\AbsEqAmp)^2=(\AbsScaleAmp\,\AbsEqAmp)^2.
\]
These are exactly \eqref{eq:rawabsA}--\eqref{eq:rawabsB}. Because \(\OriginMap\) is bijective onto the frozen signed amplitude sector, these are the same raw-lift formulas already recorded there.
\end{proof}

\begin{definition}[Induced orbit-level coefficients]
\label{def:isocoeff}
The induced orbit-level isotropic coefficient is
\begin{equation}
\RespShapeIso
:=
\frac{\RespShapeCov\,\RespShapeRelax}{\RespShapeCov+\RespShapeRelax}
=
\frac{(\AbsCovAmp)^2(\AbsRelaxAmp)^2}{(\AbsCovAmp)^2+(\AbsRelaxAmp)^2}.
\label{eq:iso}
\end{equation}
\end{definition}

\begin{theorem}[Same phase-neutral minimizing branch]
\label{thm:branch}
The orbit-level energy induced by Definition \ref{def:deepmap} is the same frozen-interface energy as in the setup, raw-orbit, and amplitude-lift notes, now with coefficients \eqref{eq:deepshapeA}--\eqref{eq:deepshapeB}. Its distinguished minimizing constant branch is therefore
\begin{equation}
\ScaleMod(x)\equiv (\AbsScaleAmp)^2,
\qquad
U(x)\equiv (\AbsScaleAmp\,\AbsEqAmp)^2,
\qquad
V(x)\equiv 0,
\qquad
\RespVacPhase(x)\equiv\phi_0,
\label{eq:branch}
\end{equation}
for arbitrary constant \(\phi_0\in\mathbb R/2\pi\mathbb Z\). In particular, the minimizing branch remains phase-neutral.
\end{theorem}

\begin{proof}
By Definition \ref{def:deepmap}, the deeper sector produces exactly the same orbit-shape coefficients and positive modulus as the amplitude-lift note, only rewritten in terms of absolute amplitudes rather than signed representatives. By Theorem \ref{thm:rawlift}, the raw-orbit lift is also unchanged. The minimizing-family computation in the frozen orbit-level energy therefore gives exactly the same constant branch as before, which becomes \eqref{eq:branch} after substituting the present absolute-amplitude formulas. Because \(\phi_0\) remains arbitrary and constant while \(V\equiv0\), the branch is still phase-neutral.
\end{proof}

\begin{theorem}[Same downstream density/current block, primitive bridge-order block, canonical profile, and dressed bridge map]
\label{thm:downstream}
For the deeper sector \(\OriginSector\), the downstream package \(\DownPkg\) is recovered unchanged. In particular:
\begin{enumerate}
    \item the same downstream density/current block is recovered;
    \item the same primitive bridge-order block is recovered;
    \item the same phase-neutral minimizing branch of Theorem \ref{thm:branch} is recovered;
    \item the same canonical profile is recovered, now written as
    \begin{equation}
    \CurvProfileCan(x)
    =
    1+\RespShapeEq\,x^2
    =
    1+(\AbsEqAmp)^2x^2
    =
    1+\frac{\RespRawEq}{\RespRawRef}x^2;
    \label{eq:profileabs}
    \end{equation}
    \item the same dressed bridge map is recovered.
\end{enumerate}
\end{theorem}

\begin{proof}
Definition \ref{def:deepmap} shows that the present deeper sector feeds the frozen interface only through the induced pair \((\ShapeTuple,\ScaleMod)\). Theorem \ref{thm:rawlift} shows that the same raw tuple is then generated. Theorem \ref{thm:branch} shows that the same phase-neutral minimizing branch survives. By Definition \ref{def:downstream}, all recorded downstream constructions depend on deeper data only through exactly those objects. Therefore every component of \(\DownPkg\) remains unchanged.

The canonical profile formula \eqref{eq:profileabs} is immediate from \(\RespShapeEq=(\AbsEqAmp)^2\) and \(\RespRawEq/\RespRawRef=\RespShapeEq\).
\end{proof}

\begin{remark}
Theorem \ref{thm:downstream} is the exact reason the present note is honest. It deepens the origin of the amplitude layer itself without trying to earn novelty by rewriting formulas that are already fixed positively downstream.
\end{remark}

\section{What Changes and What Does Not}

\begin{definition}[Polarity-forgetting relation]
\label{def:polquot}
Write \(m\in\AbsAmpSector\) for an absolute-amplitude block and \(p,p'\in\PolaritySector\) for two polarity blocks. Define an equivalence relation \(\PolarityQuot\) on \(\OriginSector\) by
\begin{equation}
(m,p)\PolarityQuot(m,p')
\label{eq:polquot}
\end{equation}
for arbitrary \(m\in\AbsAmpSector\) and arbitrary \(p,p'\in\PolaritySector\).
\end{definition}

\begin{proposition}[The polarity-forgetting quotient is exactly the old sign quotient]
\label{prop:samequot}
Under the bijection \(\OriginMap\) of Proposition \ref{prop:bijection}, the polarity-forgetting relation \(\PolarityQuot\) on \(\OriginSector\) corresponds exactly to the independent sign relation \(\SignQuot\) on the signed amplitude sector \(\LiftSector\). Consequently,
\begin{equation}
\OriginSector/\PolarityQuot
\;\cong\;
\LiftSector/\SignQuot
\;\cong\;
\AbsAmpSector.
\label{eq:quotientchain}
\end{equation}
\end{proposition}

\begin{proof}
By \eqref{eq:originmapexplicit}, changing only the polarity labels while holding the positive absolute amplitudes fixed changes only the signs of the six signed amplitudes. This is precisely the sign action defining \(\SignQuot\) in the amplitude-lift note. Conversely, every independent sign change in \(\LiftSector\) is realized uniquely by changing the corresponding polarity labels in \(\OriginSector\) while leaving the absolute amplitudes fixed. Therefore the two quotient relations match under \(\OriginMap\).

The quotient space is canonically identified with \(\AbsAmpSector\) because after forgetting polarity labels the only surviving data are the positive absolute amplitudes.
\end{proof}

\begin{definition}[Verdict-changing deeper structure]
\label{def:verdictchange}
At the present stage, a deeper continuation changes the current scale-orbit or residual vacuum-angle verdicts only if it provides at least one of the following:
\begin{enumerate}
    \item a quotient identifying distinct positive values of \(\AbsScaleAmp\), or equivalently distinct positive moduli \(\ScaleMod=(\AbsScaleAmp)^2\);
    \item a locking mechanism or true quotient fixing or identifying distinct constant vacuum angles \(\RespVacPhase=\phi_0\).
\end{enumerate}
\end{definition}

\begin{theorem}[The present deeper sector introduces no new verdict-changing structure]
\label{thm:verdict}
The oriented absolute-amplitude sector \(\OriginSector\) introduces no verdict-changing structure in the sense of Definition \ref{def:verdictchange}. Therefore:
\begin{enumerate}
    \item the positive scale orbit remains a canonical-normalization orbit rather than a proved redundancy;
    \item the global \(O(2)\) vacuum angle remains residual.
\end{enumerate}
\end{theorem}

\begin{proof}
By Proposition \ref{prop:samequot}, the only quotient in the present note is the polarity-forgetting quotient \(\PolarityQuot\), which is exactly the old finite sign quotient rewritten one layer deeper. It forgets only discrete polarity labels while leaving every positive absolute amplitude, including \(\AbsScaleAmp\), unchanged. Distinct positive values of \(\AbsScaleAmp\) and hence distinct positive moduli \(\ScaleMod=(\AbsScaleAmp)^2\) are therefore not identified. This gives item (1).

Likewise, the present deeper sector introduces no new angular variable, no anisotropic locking term, and no quotient on constant values of \(\phi_0\). The minimizing branch of Theorem \ref{thm:branch} still carries arbitrary constant \(\phi_0\in\mathbb R/2\pi\mathbb Z\). Therefore the global \(O(2)\) vacuum angle remains residual, giving item (2).
\end{proof}

\begin{corollary}[When a later verdict revision would be honest]
\label{cor:revision}
The current scale-orbit and residual-angle verdicts may be revised only if some later still-deeper sector supplies a genuine positive-modulus quotient or a genuine vacuum-angle locking or angle quotient beyond the polarity-forgetting structure of the present note.
\end{corollary}

\begin{proof}
This is the contrapositive reading of Theorem \ref{thm:verdict}. Without genuinely new quotient or locking data, the present verdicts remain.
\end{proof}

\begin{remark}
Corollary \ref{cor:revision} is intentionally strict. The present note \emph{does} make the old finite sign quotient more transparent by tracing it to explicit polarity labels, but that is not a new positive-modulus quotient and it is not vacuum-angle locking.
\end{remark}

\section{Claim Boundary}

This note does not claim that the oriented absolute-amplitude sector \(\OriginSector\) is the unique final substrate completion below the signed amplitude layer. It does not claim that the polarity labels are already dynamical observables, and it does not claim that every conceivable still-deeper completion must factor through the particular positive/polarity split written here. It also does not claim that the positive scale orbit or the global \(O(2)\) vacuum angle have already been removed by a deeper redundancy or locking mechanism.

Its claim is narrower. The six signed amplitudes feeding the primitive-class amplitude-lift note are no longer primitive either. They are generated exactly by the explicit origin map \(\OriginMap\) from a deeper oriented absolute-amplitude sector. Through the already-frozen amplitude lift and setup map, that deeper sector reproduces the same readout block, the same ordered kinetic pair, the same radial anchor pair, the same positive orbit modulus, the same raw-orbit lift, and the same downstream package. The only quotient present is the finite polarity-forgetting quotient corresponding exactly to the old sign quotient of signed-amplitude representatives, and that quotient does not change the current scale-orbit or residual-angle verdicts.

\section{Conclusion}

The next honest manuscript step below the primitive-class amplitude-lift note was to stop leaving the six signed amplitudes themselves as the deepest recorded layer. This note now does that in the least-drift way. The signed amplitudes are generated from a still-deeper oriented absolute-amplitude sector consisting of six positive absolute amplitudes together with six discrete polarity labels.

The positive result is exact and limited. The origin map \(\OriginMap\) is a bijection onto the signed amplitude sector, so each amplitude factors uniquely into positive magnitude and polarity. Composed with the frozen amplitude lift and the frozen setup map, this deeper sector reproduces the same readout block, the same ordered kinetic pair, the same radial anchor pair, the same positive orbit modulus, the same raw-orbit lift, the same density/current block, the same primitive bridge-order block, the same phase-neutral minimizing branch, the same canonical profile, and the same dressed bridge map.

The scope is equally clear. The only quotient carried by the present note is the finite polarity-forgetting quotient corresponding exactly to the old sign quotient of signed representatives. It does not identify distinct positive moduli and it does not lock or quotient the global \(O(2)\) vacuum angle. Therefore the disciplined current verdict remains in force: the positive scale orbit is still a canonical-normalization orbit rather than a proved deeper redundancy, and the vacuum angle remains residual. The next honest open burden, if the derivational line continues, is therefore deeper again: a still more primitive origin, anchoring, or selection principle for the positive absolute-amplitude block and its discrete polarity data themselves, with the verdicts revisited there only if genuinely new quotient or locking structure appears.

\end{document}
